% ===================================================================
%  TAFEL Nr. 8 — Die Ernte. --- Jagd, Weinlese, Fischfang.
%  Traduction 2026 d'apres texto/io, controlee sur texto/fr et
%  texto/en. Consigne : voir texto/de/10-tafel-01.tex.
%
%  L'appel de note est dans le TITRE de la scene — « Die Ernte (*) » —
%  et non dans un alinea : c'est ainsi que l'ido le compose, et le
%  relieur d'appels accepte les blocs d'apparat pour cette raison.
%
%  TROIS DATES SUR UNE SEULE PAGE, ET CHACUNE SE LIT DANS LA FORME
%  DES MOTS. C'est la page la plus profonde de la colonne.
%
%  1. AVANT ROME — LES GRAINS. Tout ce qui touche au ble de la
%     premiere scene est germanique et n'a jamais bouge : Ernte,
%     Sense, Sichel, Garbe, Halm, Ähre, Korn, Stroh, Dreschflegel,
%     Schober, Scheune, Mühle, Mahd. Ce sont les plus vieux mots de
%     toute la colonne, plus vieux que l'ecriture allemande.
%
%  2. ROME — LA VIGNE. Tout ce qui touche au vin de la seconde scene
%     est LATIN, sans une exception :
%        Wein      < vinum        Winzer   < vinitor
%        Weinberg  (le mont, allemand ; la vigne, latine)
%        Rebe      < ripa/repere  Most     < mustum
%        Kelter    < calcatura    Trichter < tractarius
%        Essig     < acetum       Kelch    < calix
%        Bütte     < buttis       Fass     (germanique, mais…)
%        Böttcher  < buttis       Küfer    < cupa
%     Le tonnelier de l'alinea 4 s'appelle « Böttcher » au nord et
%     « Küfer » au sud, et LES DEUX MOTS SONT LATINS : le tonneau est
%     romain ou qu'on se tienne. La vigne est montee le Rhin et la
%     Moselle avec la legion, et elle a garde son vocabulaire entier.
%
%  3. ET LA BIERE, DEUX PHRASES PLUS LOIN, EST GERMANIQUE ENTIERE :
%     Bier, Hopfen, Malz, Gerste, Brauer, Sud, Bottich, Darre. L'ido
%     de l'alinea 6 met les deux cote a cote — « dans certains pays,
%     ou la vigne est tres rare, on cultive le houblon pour faire de
%     la biere » — et l'allemand est justement le pays qui est des
%     deux cotes de cette phrase. LA FRONTIERE DU VIN ET DE LA BIERE
%     PASSE DANS LE DICTIONNAIRE ALLEMAND, et c'est la frontiere de
%     l'empire romain.
%
%  4. 1880 — LES MACHINES. La Mähmaschine de l'alinea 10, la
%     Dreschmaschine et la Lokomobile de l'alinea 11
%     sont des composes de l'age industriel, faits sur le modele du
%     tableau 1 : deux racines, le determinant devant. Deux d'entre
%     eux sont entierement allemands ; le troisieme, la locomobile,
%     est reste latin parce qu'il est venu d'Angleterre par la France.
%
%  QUATRE COUCHES, ET ELLES SE SUIVENT DANS L'ORDRE OU LES CHOSES
%  SONT ARRIVEES. Le tableau 5 avait date les metiers par la forme de
%  leur nom ; cette page date une civilisation entiere de la meme
%  facon, et sur deux mille ans.
% ===================================================================

%%K t08-noto-1 noto
\VUnotes{143.2mm}{%
\textit{(*) Weil dieses Wort ungenau ist: ist es eine Flachsernte, eine Weinlese,
eine Apfelernte? Vielleicht wäre „}
\VUgras{mesono} \textit{“ besser, vorgeschlagen in} Mondo \textit{XIV, S. 242.
Bisher ist diese neue Wurzel von der Akademie aber nicht übernommen worden und
wartet auf die Beratung nach den üblichen Regeln der Ido-Bewegung.}%
}

%%K t08-apar-1 apar
\VUsaut{5.0mm}
\VUtitre{15.0pt}{60}{TAFEL Nr. 8}
\VUsaut{3.0mm}
\VUcentre{13.2pt}{90}{\textit{Erste Szene.}}
\VUsaut{3.0mm}
\VUpk{10.2pt}{40}{Die Ernte (*).}
\VUsaut{4.0mm}
\VUpk{10.2pt}{40}{Das Gesicht des Landes.}
\VUsaut{3.0mm}

%%K t08-c1-01-1 p
1. --- Mit dem \VUgras{Sommer} hat das Land sein Gesicht wieder gewechselt. Das
Saatgut, das der Furche anvertraut wurde, ist aufgegangen; eine kleine grüne
Pflanze ist aus dem Boden gekommen und hat ihre Ähre getragen. Das Korn hat sich
gebildet, dann ist es reif geworden, und jetzt bleibt nur noch, die Ernte
einzubringen.

%%K t08-c1-02-1 p
2. --- Die \VUgras{Feldblumen} stehen in Blüte. \VUgras{Kornblumen}
\textsuperscript{(1)}, \VUgras{Mohn} \textsuperscript{(2)} und
\VUgras{Fingerhut} \textsuperscript{(3)} heben sich mit dem heiteren Kontrast ihrer
frischen \VUgras{Blüten} vom gleichmäßigen Ton der Kornfelder ab.

%%K t08-c1-03-1 p
3. --- Die Obstbäume haben ihr Laub bekommen; das Obst ist reif geworden. Der kleine
\VUgras{Kirschbaum} \textsuperscript{(4)} hängt voller schöner \VUgras{Kirschen}
\textsuperscript{(5)}, die die Kinder mit Vergnügen pflücken und essen.

%%K t08-c1-04-1 p
4. --- Die Herde kann jetzt den ganzen Tag im Freien bleiben.

%%K t08-c1-04-2 p
Die \VUgras{Lämmer} \textsuperscript{(6)} sind gewachsen und zu schönen
\VUgras{Mutterschafen} \textsuperscript{(7)} und schönen \VUgras{Widdern}
\textsuperscript{(8)} mit dichtem Vlies geworden. Sie weiden friedlich auf dem
\VUgras{Gras} der \VUgras{Weide} \textsuperscript{(9)}, die mit \VUgras{Gänseblümchen}
gesprenkelt ist.

%%K t08-c1-04-3 p
Bald werden diese Tiere wegen ihrer \VUgras{Wolle} geschoren, aus der das Tuch
unserer Kleider gemacht wird.

%%K t08-tit-2 sub
\VUsaut{5.0mm}
\VUpk{10.2pt}{40}{Der Hirt.}
\VUsaut{3.4mm}

%%K t08-c1-05-1 p
5. --- Der \VUgras{Hirt} \textsuperscript{(10)} scheint nicht viel auf sie zu achten.
Seine einzige Sorge ist das Gewitter, das am Horizont vorbeizieht, und der
\VUgras{Schäfer} \textsuperscript{(13)}, der seine \VUgras{Feldflasche}
\textsuperscript{(14)} an die Lippen hebt, wirkt noch sorgloser.

%%K t08-c1-06-1 p
6. --- Die Hitze ist drückend: sogar der \VUgras{Hund} \textsuperscript{(15)} kommt,
um sich abzukühlen; er schlappt das frische Wasser des \VUgras{Bachs}
\textsuperscript{(16)} und scheucht einen armen kleinen \VUgras{Frosch}
\textsuperscript{(17)} auf, der im Gras am Ufer gehockt hatte. Er springt in das
klare Wasser, wo die \VUgras{Seerosen} \textsuperscript{(18)} blühen.

%%K t08-c1-07-1 p
7. --- Eine \VUgras{Bachstelze} \textsuperscript{(19)} badet neben der kleinen
\VUgras{Quelle} \textsuperscript{(20)}, die zwischen zwei Felsen hervorsprudelt und
sich unter dem \VUgras{Dorngestrüpp} \textsuperscript{(21)} und den
\VUgras{Weißdornbüschen} \textsuperscript{(22)} versteckt.

%%K t08-tit-3 sub
\VUsaut{5.0mm}
\VUpk{10.2pt}{40}{Die Ernte.}
\VUsaut{3.4mm}

%%K t08-c1-08-1 p
8. --- Für die Landkinder ist die Kornernte eine sehr vergnügliche Zeit. Sie
schlagen fröhliche \VUgras{Purzelbäume} \textsuperscript{(24)} auf den
\VUgras{Strohhaufen} \textsuperscript{(23)}, sie helfen den \VUgras{Schnittern}
\textsuperscript{(26)}, und während diese das \VUgras{Kornfeld}
\textsuperscript{(27)} mit ihren \VUgras{Sensen} \textsuperscript{(25)} mähen, die am
\VUgras{Wetzstein} \textsuperscript{(30)} gut geschärft sind, sammeln sie das Korn
und binden es zu \VUgras{Garben} \textsuperscript{(31)} oder kleinen
\VUgras{Bündeln} \textsuperscript{(32)}.

%%K t08-c1-09-1 p
9. --- Manchmal dürfen die Größten mit einer \VUgras{Sichel} \textsuperscript{(12)}
die \VUgras{goldenen Halme} \textsuperscript{(28)} schneiden, die die schweren
\VUgras{Ähren} \textsuperscript{(29)} voller Körner tragen; und die Kleinen fangen,
während sie das \VUgras{Vieh} \textsuperscript{(33)} hüten, das auf der
\VUgras{Wiese} \textsuperscript{(34)} im Schatten der großen \VUgras{Linde}
\textsuperscript{(35)} weidet, die schönen \VUgras{Schmetterlinge} mit ihrem
\VUgras{Netz} \textsuperscript{(36)}.

%%K t08-c1-09-2 p
Manche klettern sogar auf den \VUgras{Wagen} \textsuperscript{(37)}, um die Garben zu
ordnen, die ihnen auf einer \VUgras{Heugabel} \textsuperscript{(38)} hinaufgereicht
werden.

%%K t08-c1-10-1 p
10. --- Auf der ganzen \VUgras{Ebene} \textsuperscript{(39)}, wo die \VUgras{Lerchen}
\textsuperscript{(41)} auffliegen, herrscht ein großes Treiben. Alle sind mit der
Ernte beschäftigt. Aber während die Armen weiter das alte Werkzeug gebrauchen ---
\VUgras{Sensen}, \VUgras{Sicheln} und \VUgras{Dreschflegel}
\textsuperscript{(40)} ---, lassen die reichen Gutsbesitzer ihren Weizen oder ihren
\VUgras{Roggen} \textsuperscript{(43)} mit einer \VUgras{Mähmaschine}
\textsuperscript{(42)} schneiden. Dann werden große \VUgras{Schober}
\textsuperscript{(44)} an Ort und Stelle aufgesetzt, ohne dass man die Garben in die
Scheune fahren müsste.

%%K t08-c1-11-1 p
11. --- Dann wird eine \VUgras{Dreschmaschine} \textsuperscript{(47)} herangefahren,
die eine \VUgras{Lokomobile} \textsuperscript{(45)} über einen
\VUgras{Treibriemen} \textsuperscript{(46)} antreibt, und so wird das Korn vom
\VUgras{Stroh} getrennt, ohne das Feld zu verlassen, auf dem es gesät wurde.

%%K t08-tit-4 sub
\VUsaut{5.0mm}
\VUpk{10.2pt}{40}{Das Gewitter.}
\VUsaut{3.4mm}

%%K t08-c1-12-1 p
12. --- Zur Erntezeit brechen oft heftige \VUgras{Gewitter} \textsuperscript{(53)}
los. Zuerst hört man von fern das Grollen des \VUgras{Donners}; ein starker
\VUgras{Westwind} \textsuperscript{(54)} kommt auf und biegt die größten Bäume des
\VUgras{Waldes} \textsuperscript{(49)}, bis sie ächzen. Schwarze \VUgras{Wolken}
\textsuperscript{(56)} ziehen über den Himmel; sie werden von den blendenden
Zickzacken des \VUgras{Blitzes} \textsuperscript{(55)} durchschnitten. Der
\VUgras{Regen}, und manchmal der \VUgras{Hagel}, fängt an zu fallen und drückt die
Ernte nieder, die geschnitten werden soll.

%%K t08-c1-13-1 p
13. --- Oft setzt der Blitz ein Gebäude in Brand. Im vorigen Jahr zum Beispiel ist
das Dach der \VUgras{Windmühle} \textsuperscript{(50)} zu Asche verbrannt, und zwei
ihrer \VUgras{Flügel} \textsuperscript{(51)} sind zerschlagen worden.

%%K t08-c1-14-1 p
14. --- Nach ein paar Stunden kehrt die Ruhe zurück. Ein leuchtender
\VUgras{Regenbogen} \textsuperscript{(52)} erscheint am Horizont, und die Natur nimmt
das Leben wieder auf, das eine Weile unterbrochen war. Die Landleute, die sich in
den \VUgras{Hütten} \textsuperscript{(48)} untergestellt hatten, gehen wieder an die
Arbeit und machen sich tapfer daran, den erlittenen Schaden zu beheben.

%%K t08-tit-5 sub
\VUsaut{6.0mm}
\VUcentre{13.2pt}{90}{\textit{Zweite Szene.}}
\VUsaut{3.6mm}

%%K t08-tit-6 sub
\VUpk{10.2pt}{40}{Jagd. --- Weinlese.}
\VUsaut{1.4mm}
\VUpk{10.2pt}{40}{Fischfang.}
\VUsaut{4.0mm}

%%K t08-c2-01-1 p
1. --- Gegen Ende \VUgras{August} oder Mitte \VUgras{September}, je nach Gegend,
beginnt die Jagdzeit. Kaum haben die ersten Sonnenstrahlen die \VUgras{Eidechsen}
\textsuperscript{(2)} aus ihren Löchern gelockt, da zieht der \VUgras{Jäger}
\textsuperscript{(5)} schon mit seinen \VUgras{Hunden} \textsuperscript{(4)} los.

%%K t08-c2-02-1 p
2. --- Er hat seinen derben \VUgras{Jagdanzug} \textsuperscript{(9)} aus dickem Tuch
angezogen und lederne \VUgras{Gamaschen}, mit denen er durch das nasse Gras gehen
kann, wo sich die \VUgras{Kröten} \textsuperscript{(1)} verstecken; er hat sein
\VUgras{Gewehr} \textsuperscript{(6)} und seine \VUgras{Jagdtasche}
\textsuperscript{(10)} genommen, und schon ist er unterwegs.

%%K t08-c2-03-1 p
3. --- Der Eifer der Jagd führt ihn mitten in einen Weinberg, wo die Lese in vollem
Gang ist. Die Kinder des Winzers, die sonst die Ziegen am Wegrand hüten, lassen sie
heute weiden, wo sie wollen, und nachdem sie einen alten \VUgras{Ziegenbock}
\textsuperscript{(3)} an einen Pfahl gebunden haben, der seine Freiheit sicher zum
Weglaufen nutzen würde, haben auch sie ihren Anteil am Vergnügen. Der eine nimmt
eine Traube aus einer \VUgras{Lesebütte} \textsuperscript{(12)} und drückt zum
Vergnügen den \VUgras{Saft} \textsuperscript{(14)} in einen \VUgras{Becher}
\textsuperscript{(13)}, um Most zu machen (ungegorenen Wein). Ein anderer setzt sich
einen \VUgras{Laubkranz} \textsuperscript{(15)} auf, den er eben geflochten
hat. Neben ihnen kommen dreiste \VUgras{Spatzen} \textsuperscript{(16)} bis an die
Kelter, um Trauben zu stehlen.

%%K t08-c2-04-1 p
4. --- Während die Kinder sich vergnügen, arbeiten die Männer. Die
\VUgras{Leser} \textsuperscript{(18)} tragen die Trauben in \VUgras{Rückenbütten}
\textsuperscript{(17)} in den Keller hinunter; ein \VUgras{Böttcher}
\textsuperscript{(19)} bessert die \VUgras{Fässer} \textsuperscript{(20)} aus, prüft,
ob die \VUgras{Dauben} \textsuperscript{(22)} und die \VUgras{Böden}
\textsuperscript{(21)} heil sind, und ersetzt die zerbrochenen \VUgras{Reifen}
\textsuperscript{(23)}. Er ist es auch, der die großen \VUgras{Bottiche}
\textsuperscript{(25)} gemacht hat, in die die \VUgras{Trauben}
\textsuperscript{(26)} zum Gären geschüttet werden.

%%K t08-c2-05-1 p
5. --- Wenn die Gärung vorbei ist, wird der Wein abgezogen und der Rest auf die
\VUgras{Kelter} \textsuperscript{(28)} gebracht; indem man die große
\VUgras{Spindel} \textsuperscript{(29)} nach und nach anzieht, presst man den Saft
heraus, der in einen \VUgras{Zuber} \textsuperscript{(32)} läuft und noch mehr
\VUgras{Wein} \textsuperscript{(31)} gibt.

%%K t08-tit-8 sub
\VUsaut{6.0mm}
\VUpk{10.2pt}{40}{Bier und Apfelwein.}
\VUsaut{3.4mm}

%%K t08-c2-06-1 p
6. --- An der Mauer des \VUgras{Kellers} \textsuperscript{(27)} rankt ein Trieb
\VUgras{Hopfen} \textsuperscript{(30)} empor. Hier ist er nur eine Zierpflanze, aber
in manchen Ländern, wo die Rebe sehr selten ist, baut man Hopfen an, um
\VUgras{Bier} zu machen.

%%K t08-c2-07-1 p
7. --- Der kleine \VUgras{Apfelbaum} \textsuperscript{(33)} hängt voller Äpfel, aus
denen ausgezeichneter \VUgras{Apfelwein} wird, wenn sie reif sind. In ihrem
\VUgras{Käfig} \textsuperscript{(34)} sehen der \VUgras{Kanarienvogel}
\textsuperscript{(35)} und der \VUgras{Stieglitz} \textsuperscript{(36)} den Spatzen
neidisch zu, die in Freiheit herumhüpfen.

%%K t08-c2-08-1 p
8. --- So weit das Auge reicht, stehen am Hang der Hügel Reihen von
\VUgras{Weinpfählen} \textsuperscript{(40)}, die die prächtigen \VUgras{Rebstöcke}
\textsuperscript{(39)} stützen, schwer von \VUgras{Trauben}.

%%K t08-c2-08-2 p
Das ganze Jahr über hat der \VUgras{Winzer} \textsuperscript{(42)} gearbeitet, um
seinen \VUgras{Weinberg} \textsuperscript{(38)} zu pflegen, und jetzt wird er mit
einer guten Ernte für seine Mühe belohnt. Die \VUgras{Leserinnen}
\textsuperscript{(41)} füllen die Bottiche auf dem Wagen, den \VUgras{Maultiere}
\textsuperscript{(43)} ziehen, und alle sind fröhlich.

%%K t08-tit-9 sub
\VUsaut{5.0mm}
\VUpk{10.2pt}{40}{Fischfang.}
\VUsaut{3.4mm}

%%K t08-c2-09-1 p
9. --- Dasselbe gilt für die \VUgras{Angler} \textsuperscript{(44)}, die bei
Tagesanbruch gekommen sind und sich mit ihren \VUgras{Angelruten}
\textsuperscript{(45)} ans Wasser gesetzt haben.

%%K t08-c2-09-2 p
Die \VUgras{Fische} \textsuperscript{(47)} beißen an, sobald sie ihre
\VUgras{Schnur} \textsuperscript{(46)} ins Wasser lassen, und sie haben kaum Zeit,
den \VUgras{Köder} zu erneuern, bevor ein neues Opfer am Ende der Schnur zappelt.

%%K t08-c2-09-3 p
Der \VUgras{Netzfischer} \textsuperscript{(48)} allerdings, der sein Netz von seinem
\VUgras{Kahn} \textsuperscript{(49)} aus mitten in den \VUgras{Fluss}
\textsuperscript{(50)} wirft, fängt weit mehr.

%%K t08-c2-10-1 p
10. --- Der Herbst ist die Zeit der schönen Ausflüge: zu Fuß, zu \VUgras{Pferd}
\textsuperscript{(52)}, mit dem \VUgras{Fahrrad} \textsuperscript{(53)}, mit dem
\VUgras{Auto} \textsuperscript{(54)}. Die \VUgras{Straßen} \textsuperscript{(51)} sind
sauber, und man nutzt die letzten schönen Tage, weil die \VUgras{Schwalben}
\textsuperscript{(56)} sich schon auf den Dächern sammeln, um in wärmere Länder zu
ziehen. Die Blätter des alten \VUgras{Nussbaums} \textsuperscript{(57)} werden gelb,
die \VUgras{Nüsse} fallen, und die hohen \VUgras{Bäume}, die in einer
\VUgras{Gruppe} \textsuperscript{(58)} auf der \VUgras{Anhöhe}
\textsuperscript{(59)} stehen, werden bald kahl sein.

%%K t08-c2-11-1 p
11. --- Die \VUgras{Sonne} \textsuperscript{(60)} geht später auf, die Tage werden
kürzer, am \VUgras{Morgen} spürt man schon eine ziemlich scharfe Kühle, und Flüge
von \VUgras{Schnepfen} \textsuperscript{(61)} verlassen bereits unser unwirtliches
Klima.
\VUsaut{6.0mm}
\VUfilet{22mm}
